% GENERATED FILE. Edit canonical lessons, then run npm run generate:books.
% canonical-source-hash: fnv1a64:07b708a11e0d6bc0
% canonical-lessons: FR-C432-gerant, FR-C432-desagrement, FR-C432-signaler, FR-C432-menage, FR-C432-defectueux

\chapter{Parler au gérant}
\label{ch:fr-parler-au-gerant}
\hlchaptermodality{432}

\begin{chapteropening}
\textbf{By the end of this chapter:} \emph{I can take a problem to the manager: the manager, the inconvenience, reporting a problem, the cleaning and something faulty.}

Complete the last lesson of chapter 432: I can take a problem to the manager: the manager, the inconvenience, reporting a problem, the cleaning and something faulty.
\end{chapteropening}

\section[le gérant]{le gérant — the manager}

\label{lesson:FR-C432-gerant}



\begin{quote}
Before the new one: say the French for to disturb, then the French for unhappy.
\end{quote}

\subsection*{You'll want to know: le gérant}
\textbf{le gérant} — \textquotedblleft{}the manager\textquotedblright{}.

\begin{quote}
Le gérant de l'hôtel est venu voir la chambre.
\end{quote}

\begin{quote}
\emph{The hotel manager came to check the room.}
\end{quote}

Feminine \emph{la gérante}.

\begin{cousinweb}
\emph{gérant} is built on \emph{gérer} 'to manage', a learned borrowing from Latin \emph{gerere} 'to carry on, to run' (Lewis \& Short). EtymDB lists it without a source.
\end{cousinweb}

\begin{grammarlens}[title={What you've built}]
One more word for telling, explaining and describing at B1.
\end{grammarlens}

\subsection*{Your turn}
\begin{itemize}
\raggedright
  \item \emph{Say it:} \emph{le gérant}
  \item \emph{Say it:} \emph{le gérant}, once more
  \item \emph{Say it:} a sentence of your own with \emph{le gérant}
  \item \emph{Recall:} say the French for to disturb, then the French for unhappy, then say \emph{le gérant} again
\end{itemize}

\begin{tcolorbox}[breakable,colback=teal!4,colframe=teal!35!black,title={Before you move on}]
What does le gérant mean? (\textquotedblleft{}The manager\textquotedblright{}.) Say it once more.
\end{tcolorbox}
\section[le désagrément]{le désagrément — the inconvenience, the nuisance}

\label{lesson:FR-C432-desagrement}



\begin{quote}
Before the new one: say the French for unhappy, then the French for the manager.
\end{quote}

\subsection*{You'll want to know: le désagrément}
\textbf{le désagrément} — \textquotedblleft{}the inconvenience, the nuisance\textquotedblright{}.

\begin{quote}
Nous sommes désolés pour ce désagrément.
\end{quote}

\begin{quote}
\emph{We are sorry for the inconvenience.}
\end{quote}

A polite word in apologies.

\begin{cousinweb}
\emph{désagrément} is built from \emph{dés-} and \emph{agrément} 'pleasure', from \emph{agréer} 'to please', from \emph{gré} 'liking', which goes back to Latin \emph{grātum} 'something pleasing' (Lewis \& Short). EtymDB does not list it.
\end{cousinweb}

\begin{grammarlens}[title={What you've built}]
One more word for telling, explaining and describing at B1.
\end{grammarlens}

\subsection*{Your turn}
\begin{itemize}
\raggedright
  \item \emph{Say it:} \emph{le désagrément}
  \item \emph{Say it:} \emph{le désagrément}, once more
  \item \emph{Say it:} a sentence of your own with \emph{le désagrément}
  \item \emph{Recall:} say the French for unhappy, then the French for the manager, then say \emph{le désagrément} again
\end{itemize}

\begin{tcolorbox}[breakable,colback=teal!4,colframe=teal!35!black,title={Before you move on}]
What does le désagrément mean? (\textquotedblleft{}The inconvenience, the nuisance\textquotedblright{}.) Say it once more.
\end{tcolorbox}
\section[signaler]{signaler — to report, to point out}

\label{lesson:FR-C432-signaler}



\begin{quote}
Before the new one: say the French for the manager, then the French for the inconvenience.
\end{quote}

\subsection*{You'll want to know: signaler}
\textbf{signaler} — \textquotedblleft{}to report, to point out\textquotedblright{}.

\begin{quote}
J'ai signalé la fuite à la réception.
\end{quote}

\begin{quote}
\emph{I reported the leak to reception.}
\end{quote}

\emph{Signaler un problème} — to report a problem.

\begin{cousinweb}
\emph{signaler} is built on \emph{signal}, from Late Latin \emph{signāle}, from \emph{signum} 'mark, sign' (Lewis \& Short).
\end{cousinweb}

\begin{grammarlens}[title={What you've built}]
One more word for telling, explaining and describing at B1.
\end{grammarlens}

\subsection*{Your turn}
\begin{itemize}
\raggedright
  \item \emph{Say it:} \emph{signaler}
  \item \emph{Say it:} \emph{signaler}, once more
  \item \emph{Say it:} a sentence of your own with \emph{signaler}
  \item \emph{Recall:} say the French for the manager, then the French for the inconvenience, then say \emph{signaler} again
\end{itemize}

\begin{tcolorbox}[breakable,colback=teal!4,colframe=teal!35!black,title={Before you move on}]
What does signaler mean? (\textquotedblleft{}To report, to point out\textquotedblright{}.) Say it once more.
\end{tcolorbox}
\section[le ménage]{le ménage — the cleaning, the housework; the household}

\label{lesson:FR-C432-menage}



\begin{quote}
Before the new one: say the French for the inconvenience, then the French for to report.
\end{quote}

\subsection*{You'll want to know: le ménage}
\textbf{le ménage} — \textquotedblleft{}the cleaning, the housework; the household\textquotedblright{}.

\begin{quote}
Le ménage n'a pas été fait dans notre chambre.
\end{quote}

\begin{quote}
\emph{Our room has not been cleaned.}
\end{quote}

\emph{Faire le ménage} — to do the cleaning.

\begin{cousinweb}
\emph{ménage} comes from Old French \emph{manage} 'dwelling, household' (EtymDB), usually traced to a word built on Latin \emph{manēre} 'to stay'.
\end{cousinweb}

\begin{grammarlens}[title={What you've built}]
One more word for telling, explaining and describing at B1.
\end{grammarlens}

\subsection*{Your turn}
\begin{itemize}
\raggedright
  \item \emph{Say it:} \emph{le ménage}
  \item \emph{Say it:} \emph{le ménage}, once more
  \item \emph{Say it:} a sentence of your own with \emph{le ménage}
  \item \emph{Recall:} say the French for the inconvenience, then the French for to report, then say \emph{le ménage} again
\end{itemize}

\begin{tcolorbox}[breakable,colback=teal!4,colframe=teal!35!black,title={Before you move on}]
What does le ménage mean? (\textquotedblleft{}The cleaning, the housework; the household\textquotedblright{}.) Say it once more.
\end{tcolorbox}
\section[défectueux]{défectueux — faulty, defective}

\label{lesson:FR-C432-defectueux}



\begin{quote}
Before the new one: say the French for to report, then the French for the cleaning.
\end{quote}

\subsection*{You'll want to know: défectueux}
\textbf{défectueux} — \textquotedblleft{}faulty, defective\textquotedblright{}.

\begin{quote}
La télévision est défectueuse depuis notre arrivée.
\end{quote}

\begin{quote}
\emph{The television has been faulty since we arrived.}
\end{quote}

Feminine \emph{défectueuse}.

\begin{cousinweb}
\emph{défectueux} is a learned borrowing from Medieval Latin \emph{dēfectuōsus}, from \emph{dēfectus} 'lack', from \emph{dēficere} 'to fail', built on \emph{facere} 'to do' (Lewis \& Short).
\end{cousinweb}

\begin{grammarlens}[title={What you've built}]
One more word for telling, explaining and describing at B1.
\end{grammarlens}

\subsection*{Your turn}
\begin{itemize}
\raggedright
  \item \emph{Say it:} \emph{défectueux}
  \item \emph{Say it:} \emph{défectueux}, once more
  \item \emph{Say it:} a sentence of your own with \emph{défectueux}
  \item \emph{Recall:} say the French for to report, then the French for the cleaning, then say \emph{défectueux} again
\end{itemize}

\begin{tcolorbox}[breakable,colback=teal!4,colframe=teal!35!black,title={Before you move on}]
What does défectueux mean? (\textquotedblleft{}Faulty, defective\textquotedblright{}.) Say it once more.
\end{tcolorbox}
